\section{Abstract}

Portfolio optimization remains one of the fundamental challenges in quantitative investment management, requiring the construction of portfolios that maximize expected returns while controlling investment risk. This report presents the development and evaluation of a quantitative portfolio optimization framework using historical equity data from selected companies listed on the Johannesburg Stock Exchange (JSE).

Historical daily adjusted closing prices were collected and preprocessed to compute asset returns and estimate the statistical inputs required for portfolio optimization. Asset expected returns, covariance matrices and correlation structures were estimated, with covariance shrinkage employed to improve the stability of risk estimates. Modern Portfolio Theory was subsequently applied to construct both the Minimum Variance Portfolio and the Maximum Sharpe Ratio Portfolio. The efficient frontier was generated through portfolio simulations to illustrate the trade-off between expected return and portfolio risk.

The optimized portfolio was evaluated through historical backtesting using cumulative returns, wealth growth, and maximum drawdown analysis. Portfolio performance was benchmarked against the Johannesburg Stock Exchange All Share Index (JSE ALSI) to assess the effectiveness of the optimization strategy under actual market conditions. Results indicate that the optimized portfolio delivered superior cumulative performance relative to the benchmark while maintaining an acceptable level of portfolio risk throughout the investment period.

The findings demonstrate the practical application of quantitative portfolio optimization techniques in equity portfolio construction and highlight the importance of disciplined, data-driven investment decision-making. Although the analysis assumes frictionless markets and static portfolio weights, the framework provides a solid foundation for more advanced portfolio management techniques incorporating dynamic rebalancing, transaction costs and alternative risk models.

\vspace{1cm}

\clearpage